\section{Exact vertex gauges accelerate a strict-descent disc witness}
\label{sec:epoch-gauge}

The preceding epoch policy transports one cochain through smaller sources
without changing its gauge. Reoptimizing that gauge on the new geometry
can expose a simpler coherent fibre. This continuation does so periodically,
adds exact independent coboundary replay, and obtains the same actual
nonempty-diagram disc in fewer epochs. The optimization remains optional;
source, move and final component authority are checked from the original
input rather than inferred from a span objective.

\subsection{A vertex change is an exact integral coboundary}

Let $v(t,j)$ be the actual global vertex of corner $j$ in tetrahedron $t$.
For a single integral potential $p$ on these global vertices, set
\[
 h'_{t,j}=h_{t,j}+p(v(t,j))-
                         \bigl(h_{t,0}+p(v(t,0))\bigr).
\]
The final subtraction is only row normalization. Along an oriented source
edge $a\to b$, the new edge cochain differs by $p(b)-p(a)$. A period around
any closed source cycle is consequently unchanged by telescoping. Local
triangle sums and consistency across shared edge occurrences are preserved.
The correction is integral, so this is a change of representative within
the same integral cohomology class.

\begin{proposition}[Source gauge replay is exact]
An independently prepared source, a complete vertex-potential list and the
displayed corner equation suffice to authenticate the new edge cochain.
No optimizer, claimed minimum, flow transcript or move producer is required.
\end{proposition}
\begin{proof}
Preparation reconstructs the global vertex and oriented edge classes from
actual face pairings. Every stated potential must refer to exactly one such
vertex. Checking the displayed equation at all corners verifies every
new local height difference. Checking all signed edge occurrences confirms
that those differences form the same global source cochain plus the stated
vertex coboundary. The closed-cycle period identity follows directly.
The assertion concerns the cochain class; it does not assert equality of
the represented embedded normal fibres.
\end{proof}

The native producer invokes the existing span optimizer on that current
source and cochain. It anchors each connected vertex component's irrelevant
constant, constructs the new normalized rows and independently verifies the
gauge proof before using it. It also requires nonincreasing total disc span.
These are producer checks for the chosen policy, not assumptions made by the
positive disc verifier.

\subsection{A mixed move/gauge chain retains the original source}

The local \code{cocycle-vertex-gauge-v1} proof contains the actual source
digest, sorted complete vertex/value pairs and new height rows. The checker
reconstructs the source, rejects missing/extra/Boolean vertex identities,
reads exact integers through the maintained codec, and checks every corner
and signed edge equation. Arbitrary common potential offsets cancel under
row normalization, including values requiring thousands of encoded bits.
There is no expansion proportional to such a value.

Mixed chains use \code{diagram-transport-disc-v2}. The top-level original
PD, exterior, optional shellings, initial heights, final coordinates and
terminal disc proof retain their meanings. A move step retains its original
triangulation/transport pair and full geometric replay. A gauge step has
only its source-bound gauge proof, leaves the current triangulation unchanged
and supplies the authenticated next height rows. The independent full-chain
checker accepts no gauge step under the old v1 schema. Old v1 proofs still
use their original move-only path.

Neither checker imports the gauge optimizer. Every topology-changing
Pachner step still has its own source/relative-boundary proof; every gauge
step is checked on the exact current source. A lower objective value does
not imply an essential disc, and a change in fibre topology is not described
as transport of one unchanged embedded surface. The final actual coherent
coordinates and essential-component count are independently reconstructed
and checked as before.

\subsection{Periodic use and representation-size accounting}

The epoch API exposes \code{regauge\_interval}; zero keeps the previous
policy. A positive interval invokes the optimizer after that many strict
descents, with the same global work and node allowance. Identity results
need no new chain step. The recognizer and CLI expose the interval only with
explicit epoch mode, and neither optimization nor epoch search becomes a
default. A cap or failure before a complete positive replay remains
inconclusive and reaches the complete fallback.

Let $H$ bound each normalized height-row span and $t$ the initial source
size. Its total disc span is $W=\sum_t(\max h_t-\min h_t)\le tH$.
A verified nonincreasing span optimization has each new row span at most
$W$, hence at most $tH$. In combination with the preceding $4H$ per-move
bound, $m$ geometric moves and $g$ gauge calls give height bits bounded by
\[
                 B+O\bigl(m+g\log(t+1)\bigr).
\]
Strict descent gives $m\le t(2U+1)$ and $g\le t$ for a positive interval,
so this growth is polynomial in $t,U,B$.

Proof potentials are normalized per vertex component. Along a spanning
path of its source one-skeleton, $p(b)-p(a)$ is the difference between new
and old edge cochains. There are at most $4t$ vertices, and each edge
correction is bounded by $(t+1)H$ for one gauge call. Anchoring therefore
bounds potential magnitudes by $4t(t+1)H$; their bit size is also charged.
The independent checker may accept larger harmless offsets supplied by a
caller, but reading them costs their actual encoding size.

Source preparation, the span optimizer, every gauge replay and every final
component query remain explicit costs. This bound controls representation
growth; it does not turn an unaccounted operation into a constant-time one.
The bounded geometry/restart argument and the separate essentiality and
universal reach obligations retain their preceding scope.

\subsection{Actual proof reduction and complete configured-call measurements}
All 1,521 maintained tests pass in 286.636 seconds. The 28 focused tests
pass in 19.532 seconds. They check every actual gauge edge difference as
one global vertex coboundary, source and potential forgeries, changed height
rows, missing gauge steps, v1 schema downgrade, source substitution, false
component counts and producer-disabled full mixed-chain replay. A common
$2^{9000}$ potential offset leaves all normalized rows unchanged and still
passes replay after the portable integer codec. CLI routing, shared caps,
legacy zero-interval behavior and caller exception identity also pass.

The 67.775-second audit compares the frozen ungauged epoch policy with
interval-four reoptimization on the same six actual diagrams and shared
two-million-unit allowance. Exact disabled recognizer statuses, methods and
evidence are preserved, and every zero-local-work enabled call reaches the
correct complete fallback. Fresh Regina confirms all three returned disc
witnesses. The target nonempty proof also runs through the actual enabled
recognizer and full source-bound replay.

For genus-one-miss, both epoch policies certify the unknot. The ungauged
chain takes 33 strict descents, 191 nodes and 914,561 callback work units.
The gauged chain takes 24 descents, 83 nodes and 622,768 work units, about
32 percent fewer checkpoints, reaching thirty tetrahedra instead of twenty-one.
Six optimization calls yield three nonidentity gauge changes. Its v2 proof
has four upward and twenty-eight downward geometric events plus three gauge
records, 35 steps in total. The earlier disc is not attributed to a changed
cohomology class: every intermediate correction is source-replayed as an
integral vertex coboundary.

Other controls show the cost of this policy. Optimized-positive still reaches
the work cap, after 44 descents rather than 48; eleven optimizer calls produce
two actual changes. Trefoil still reaches its bounded-stationary source after
fourteen descents; three gauge changes raise its work from 1,343,190 to
1,382,485, about three percent. Figure-eight still caps after eighteen
descents with four changes. The empty-circle and one-crossing controls
return their initial-source disc without a gauge. None of these nonpositive
outcomes is used as a knottedness certificate.

Configured-call timings use three randomized paired rounds, 48 measured
calls and sixteen warmups, including identical-baseline A/A controls. Both
policies share a one-million-unit allowance and a thousand geometric nodes.
Every source preparation, span optimization, gauge check, epoch search,
component query, full positive replay and serialized result is charged.
The preceding ungauged source and recognizer modules are frozen, rather
than disabling just one helper inside the current implementation.

\textbf{Configured native epoch queries, ungauged versus periodic gauges.}
\begin{center}
\small
\begin{tabular}{lrrrrr}
\toprule
Input & ungauged ms & gauged ms & old/new & old A/A & new A/A\\
\midrule
empty circle & 8.435 & 8.392 & 1.003 & 1.009 & 0.992\\
genus one miss & 5064.360 & 3222.290 & 1.572 & 1.001 & 1.002\\
trefoil & 5226.081 & 5066.872 & 1.031 & 0.991 & 0.999\\
\bottomrule
\end{tabular}
\end{center}

\textbf{Complete recognition using the two explicit optional epoch policies.}
\begin{center}
\small
\begin{tabular}{lrrrrr}
\toprule
Input & ungauged ms & gauged ms & old/new & old A/A & new A/A\\
\midrule
genus one miss & 4807.415 & 3024.992 & 1.589 & 1.002 & 1.002\\
\bottomrule
\end{tabular}
\end{center}


The target native query has paired old/new ratio 1.572, median 5.064 versus
3.222 seconds. Complete optional recognition has ratio 1.589, median 4.807
versus 3.025 seconds. Both arms return an independently replayed unknot
certificate, and their A/A controls are essentially at parity. The gauges,
optimizer, changed chain schema and all replay costs are included. Thus the
gain is measured completed-call work on the demonstrated geometric witness,
not an inference from a smaller objective or a capped old result.

The empty-circle control is near parity at 1.003 and reaches no gauge. The
trefoil configured native control has ratio 1.031 with both policies capped;
its A/A observations are about 0.991 and 0.999. No broad speedup or negative
knot conclusion is inferred from that row, and the complete two-million-unit
audit already records extra gauge work on its bounded-stationary source.
Ordinary recognition with the optional stage omitted still solves the small
target far faster. The optional stage, epoch mode and periodic interval
therefore retain their disabled defaults.

Runtime release \code{6fea587df}, the passing tests, six-diagram audit,
fresh Regina confirmations and paired timings are retained under
\code{data/epoch-gauge-*}.

The driver is
\code{causal\_research.epoch\_gauge}, with \code{audit} and \code{benchmark}
modes. Publication review verifies 650 runtime pins, all completed and capped
outcomes and all three original-source disc proofs with gauge optimizer,
epoch, move and search producers and source constructors disabled. The three actual gauge
records, their exact scope, the 35-step mixed chain and every strict descent
are checked explicitly. General cheap support/reach, stationary-source
recognition and quasi-polynomial unknot bounds remain unproved.

